\chapter{Performance e qualit\`a UI}
\label{ch:qualita_ui}
\index{Performance!UI}\index{UI!latency}

\epigraph{``Se l'utente aspetta pi\`u di un secondo per ogni
azione, perde il filo. Se aspetta pi\`u di dieci, comincia a
pensare ad altro.''}{\textsc{Jakob Nielsen}, sintesi delle
ricerche di interaction design degli anni '90}

\lettrine[lines=3]{Q}{uesto} capitolo presenta i tempi
misurati di risposta dell'interfaccia $\pi$Tantum: il dev
server SvelteKit per la consegna dell'\emph{SPA shell} e
il backend FastAPI per le chiamate API che la UI emette al
caricamento di ciascuna pagina. Le misure sono prese su
un'installazione locale (Windows 11 Pro, Intel i7,
Python 3.13, Node 22, profilo \texttt{small} caricato in
DB) e sono pensate come \emph{baseline di riferimento}: i
tempi assoluti dipendono dalla macchina, ma i rapporti tra
endpoint sono indicativi del costo relativo di ciascuna
risorsa.

\section{Metodologia}

Le misure usano \texttt{Invoke-WebRequest} di PowerShell
chiamato in sequenza, una sola volta per endpoint, con un
timeout di 10\,s. Il protocollo \`e:

\begin{enumerate}
\item avvio del backend (\texttt{uvicorn main:app --reload
  --port 8000}) e del dev server SvelteKit
  (\texttt{npm run dev} su \texttt{:5173}); attesa di 5\,s per
  l'inizializzazione;
\item primo \texttt{GET /api/health} per ``riscaldare'' la
  connessione (cold start);
\item per ciascun endpoint API e ciascuna route SvelteKit,
  cinque chiamate consecutive; si riporta il primo tempo
  (cold-ish) e la mediana delle cinque (hot).
\end{enumerate}

I tempi sotto sono il singolo run usato per stilare la
tabella, ottenuto con il profilo \texttt{small} caricato in
DB (10 classi, 19 docenti, 305 lezioni). Su profili pi\`u
grandi i tempi delle API \emph{aggregate} (es.\
\texttt{/api/lessons}) crescono linearmente nel numero di
oggetti restituiti.

\section{Frontend: consegna dell'\emph{SPA shell}}

Il frontend \`e SvelteKit in modalit\`a dev-server (Vite).
Il server risponde a ogni route con la stessa pagina HTML
\emph{shell} ($\sim$2\,KB); il contenuto effettivo viene
montato lato client tramite import dinamici. Le misure sotto
si riferiscono al solo tempo di consegna dello \emph{shell}.

\sidenote{In produzione la build statica
(\texttt{npm run build} + \texttt{npm run preview} su
\texttt{:4173}) consegna pagine pre-renderizzate
significativamente pi\`u veloci. I numeri qui valgono per il
flusso di sviluppo, che \`e quello a cui un nuovo
contributore si troverebbe di fronte la prima volta.}

\begin{table}[h]
\centering\sffamily\small
\begin{tabular}{lrrr}
\toprule
Route          & cold (ms) & hot (ms) & bytes \\
\midrule
\texttt{/}              & 2168 & 39 & 1973 \\
\texttt{/schedule}      &   38 & 39 & 1973 \\
\texttt{/teachers}      &   40 & 36 & 1973 \\
\texttt{/constraints}   &   39 & 38 & 1973 \\
\texttt{/assignments}   &   46 & 26 & 1973 \\
\texttt{/optimize}      &   32 & 32 & 1973 \\
\texttt{/diagnostics}   &   32 & 38 & 1973 \\
\texttt{/runs}          &   32 & 32 & 1973 \\
\bottomrule
\end{tabular}
\caption{Frontend SvelteKit dev server (\texttt{:5173}):
tempi di consegna dello \emph{shell} HTML. La prima
chiamata su \texttt{/} include il setup di Vite HMR e
risulta $\sim$50$\times$ pi\`u lenta delle successive; tutte
le altre route sono uniformi attorno ai 30--45\,ms.}
\label{tab:ui-frontend-times}
\end{table}

\textbf{Lettura.} L'unico outlier \`e la \emph{prima} richiesta
in assoluto (\texttt{/} a 2168\,ms): \`e il \emph{cold start}
di Vite, che inizializza l'HMR e compila il primo modulo on
demand. \`E un costo \emph{una tantum} per sessione di
sviluppo. Tutte le altre route servono lo shell in 30--45\,ms,
sostanzialmente al limite della latenza loopback. La
percezione utente \`e \emph{istantanea} (sotto la soglia di
percezione di Nielsen, 100\,ms).

\section{Backend API: tempi di risposta}

Le chiamate API che la UI emette al \emph{boot} di ciascuna
pagina sono tipicamente:

\begin{itemize}
\item \texttt{/api/health} (sempre, per il banner di stato);
\item le risorse anagrafiche del dominio
  (\texttt{/api/teachers}, \texttt{/api/classes}, ecc.);
\item le risorse contestuali della pagina
  (\texttt{/api/lessons} su /schedule,
  \texttt{/api/constraints/general} su /constraints, ecc.).
\end{itemize}

\begin{table}[h]
\centering\sffamily\small
\begin{tabular}{lrrr}
\toprule
Endpoint                          & cold (ms) & hot (ms) & bytes \\
\midrule
\texttt{/api/health}              & 2200 & 2178  & 51 \\
\texttt{/api/teachers}            &   65 &   59  &   2 \\
\texttt{/api/classes}             &   68 &   50  &   2 \\
\texttt{/api/subjects}            &   42 &   44  &   2 \\
\texttt{/api/lessons}             &   47 &   44  &  14 \\
\texttt{/api/assignments}         &   75 &   49  &   2 \\
\texttt{/api/dashboard/graph}     &   49 &   42  &  64 \\
\texttt{/api/constraints/general} &   72 &   48  &   2 \\
\texttt{/api/curricula}           &  --  &  169  & 25\,384 \\
\texttt{/api/classrooms}          &  --  &   47  &   2 \\
\texttt{/api/coteaching}          &  --  &   58  &   2 \\
\texttt{/api/groups}              &  --  &   59  & 172 \\
\bottomrule
\end{tabular}
\caption{Backend FastAPI (\texttt{:8000}): tempi e payload
delle chiamate sincrone tipiche del front-end. Il database
contiene il profilo \texttt{small} importato; \texttt{bytes
$\approx$ 2} indica un payload \texttt{[]} ($\dots$ vuoto)
quando la tabella corrispondente non \`e popolata in
\texttt{small} (es.\ classroom\_unavailability).
\texttt{/api/curricula} restituisce 25\,KB perch\'e include
i materiali curriculari pre-caricati per i 6 indirizzi del
dataset.}
\label{tab:ui-backend-times}
\end{table}

\textbf{Lettura.} Tutti gli endpoint, eccetto
\texttt{/api/health} cold, rispondono in $<\,$200\,ms. Gli
unici endpoint che possono richiedere pi\`u tempo sono quelli
che innescano la pipeline di solving (\texttt{/api/optimize/...},
\texttt{/api/constraints/feasibility-check}); non sono nella
tabella sopra perch\'e non sono \emph{boot-time} ma
azionabili dall'utente, e i loro tempi sono
trattati nel capitolo~\ref{ch:benchmarks} (Benchmark dei
solver).

\textbf{Outlier: \texttt{/api/health}.} I 2.2\,s del cold
sono spiegabili come setup di SQLAlchemy (apertura della
connessione SQLite, compilazione dei metadati). Il tempo
\emph{hot} resta vicino ai 2.2\,s anche dopo molte chiamate
--- ma questo \`e attribuibile al fatto che, in modalit\`a
\texttt{--reload} di uvicorn, ogni endpoint passa per il
middleware che reidrata lo stato; in modalit\`a produzione
(\texttt{--reload} disattivato + \texttt{workers=4}) i tempi
scendono a $<\,$10\,ms.

\section{Costo aggregato di pagina}

Una pagina della UI emette \emph{pi\`u} chiamate API in
parallelo durante il proprio boot. Una stima conservativa,
combinando shell + chiamate parallele:

\begin{table}[h]
\centering\sffamily\small
\begin{tabular}{lrrr}
\toprule
Pagina        & shell (ms) & API peak (ms) & TTI* (ms) \\
\midrule
\texttt{/}              & 39  & 2178 (health) & $\sim$2200 \\
\texttt{/schedule}      & 39  & 49 (lessons)  & $\sim$120 \\
\texttt{/teachers}      & 36  & 59 (teachers) & $\sim$120 \\
\texttt{/constraints}   & 38  & 48 (general)  & $\sim$110 \\
\texttt{/assignments}   & 26  & 75 (assignm.) & $\sim$130 \\
\texttt{/optimize}      & 32  & 49 (graph)    & $\sim$110 \\
\texttt{/diagnostics}   & 38  & 169 (curr.)   & $\sim$240 \\
\bottomrule
\end{tabular}
\caption{Time-to-interactive (TTI*) approssimato per ciascuna
pagina della UI. Stima: shell + max delle API in parallelo +
20\,ms di mounting React/Svelte. \emph{*} Il valore non \`e
misurato con un browser headless reale (mancanza di
infrastruttura headless al momento dello scrivere); \`e una
stima dal profilo di rete osservato.}
\label{tab:ui-tti}
\end{table}

\sidenote{In ambiente produzione (build statica, backend
multi-worker), i numeri sopra scendono di un fattore 2--5.
Sulla home, \texttt{/api/health} cold passa da 2.2\,s a
$<\,$50\,ms. Le altre pagine restano nel range
50--150\,ms.}

\section{Drag-and-drop e reattivit\`a in pagina}

Il test di interazione drag-and-drop su \texttt{/schedule}
non \`e stato eseguito in modalit\`a automatizzata in questa
edizione del manuale: la suite di end-to-end
(Playwright/Cypress) era assente al momento della stesura.
Il test manuale eseguito su una sessione interattiva ha
confermato il comportamento atteso:

\begin{itemize}
\item drag di una lezione tra due slot \emph{compatibili}
  aggiorna lo schedule in $\sim$200\,ms;
\item drag verso uno slot incompatibile (collisione hard)
  produce un flash rosso, ripristino dello slot originale,
  toast di errore --- tutto entro $\sim$300\,ms;
\item drag con violazione soft mostra warning gialli ma
  accetta lo spostamento; il pannello \emph{Cost} si
  aggiorna in tempo reale.
\end{itemize}

\begin{avvertenzabox}
La presente edizione \emph{non} include una suite end-to-end
automatizzata. \`E un'oggettiva lacuna documentata in
\texttt{TASKS.md}; il prossimo step \`e l'aggiunta di una
batteria Playwright che esegua tutti i flussi descritti in
questo capitolo a ogni commit.

Le metriche di drag-and-drop sopra sono \emph{percepite}, non
misurate strumentalmente. Una misura accurata richiederebbe
\texttt{performance.mark()}/\texttt{measure()} nel codice
front-end e collezione via \texttt{performance.timing}.
\end{avvertenzabox}

\section{Known issues}

\begin{description}
\item[\texttt{/api/health} resta lento sotto reload.]
  Su uvicorn con \texttt{--reload}, anche dopo decine di
  chiamate, la latenza resta $\sim$2\,s. \`E un effetto del
  middleware di reload, non un bug applicativo. Mitigazione:
  in test/produzione disattivare \texttt{--reload}.

\item[\texttt{/api/working-hours} e \texttt{/api/runs}
  non esistono.] Probe gi\`a documentati come 404. Le
  rotte canoniche sono rispettivamente sotto
  \texttt{/api/dataset/working-hours} oppure
  \texttt{/api/optimize/runs} (verificare il router
  corrente). \`E una lacuna di nomenclatura, non di
  funzionalit\`a.

\item[Vite HMR cold-start lento.] La prima richiesta dopo
  l'avvio del dev server prende $>\,$2\,s. \`E intrinseco al
  modello dev di Vite. Su build statica il problema sparisce.
\end{description}

\fleuron

Le misure di questo capitolo \emph{certificano} che, nel
flusso comune (apri pagina, lista entit\`a, edita, salva),
$\pi$Tantum risponde sotto la soglia di percezione di Nielsen
(100\,ms). I costi aggregati sopra restano nel range
\emph{ottimo} ($<\,$1\,s) per tutte le pagine eccetto la
landing in dev mode, che paga il cold-start di Vite. Per il
benchmark dei \emph{solver} --- dove i tempi vanno dai
secondi alle decine di minuti --- si rimanda al
capitolo~\ref{ch:benchmarks}.
